% EGE 2027, task 6. Source IDs and URLs are kept in comments.
% Editable task bank. Append new BankTask blocks at the end of the document.
\documentclass[12pt,a4paper]{article}
\usepackage{physics-textbook}
\usepackage{graphicx,needspace}
\newif\ifShowSolutions
\ifdefined\TutorEdition
  \ShowSolutionstrue
\else
  \ShowSolutionsfalse
\fi
\newcounter{banktask}
\newcommand{\BankTask}[3]{%
  \refstepcounter{banktask}%
  \Needspace{35mm}\par\medskip\noindent\begin{minipage}{\linewidth}%
  \noindent\textbf{\thebanktask.}\quad #2\par
  \ifShowSolutions
    #3\par
    \ifdefined\TutorEdition
      \ifcsname TutorSource#1\endcsname
        \smallskip{\small\csname TutorSource#1\endcsname\par}%
      \fi
    \fi
  \fi
  \end{minipage}\par\medskip
}
\begin{document}


% Source EC9E02: https://ege.fipi.ru/bank/questions.php?search=1&proj=BA1F39653304A5B041B656915DC36B38&qid=EC9E02
\BankTask{EC9E02}{
На гладком горизонтальном столе брусок массой $M$, прикреплённый к вертикальной стене лёгкой горизонтальной пружиной жёсткостью $k$, совершает свободные гармонические колебания с амплитудой $A$.
\TaskMatchingColumns[.53\linewidth]{Физические величины}{
А) частота колебаний бруска\par
Б) амплитуда ускорения бруска
}{Формулы}{
1) $\frac{1}{2\pi}\sqrt{\frac{M}{k}}$\par
2) $A\frac{k}{M}$\par
3) $\frac{1}{2\pi}\sqrt{\frac{k}{M}}$\par
4) $A\sqrt{\frac{k}{M}}$\par
}\begin{center}\begin{tabular}{|c|c|}\hline А & Б \\\hline\rule{0pt}{1.2em}\hspace{1cm} & \hspace{1cm} \\\hline\end{tabular}\end{center}
}{
\smallskip\noindent\textbf{Решение от GPT}\par
$\omega=\sqrt{k/M}$, поэтому $\nu=\omega/(2\pi)$ (формула $3$), $a_{\max}=\omega^2A=Ak/M$ (формула $2$). Ответ: $32$.
}

% Source 1D0119: https://ege.fipi.ru/bank/questions.php?search=1&proj=BA1F39653304A5B041B656915DC36B38&qid=1D0119
\BankTask{1D0119}{
Шайба массой $m$, скользящая по гладкой горизонтальной поверхности со скоростью $v$, абсолютно неупруго сталкивается с покоящейся шайбой массой $M$.
\TaskMatchingColumns[.53\linewidth]{Физические величины}{
А) кинетическая энергия покоившейся шайбы после столкновения\par
Б) импульс налетающей шайбы после столкновения
}{Формулы}{
1) $\frac{m^2v}{m+M}$\par
2) $mv$\par
3) $\frac{m^2Mv^2}{2(m+M)^2}$\par
4) $\frac{Mv^2}{2}$\par
}\begin{center}\begin{tabular}{|c|c|}\hline А & Б \\\hline\rule{0pt}{1.2em}\hspace{1cm} & \hspace{1cm} \\\hline\end{tabular}\end{center}
}{
\smallskip\noindent\textbf{Решение от GPT}\par
Из сохранения импульса общая скорость $u=mv/(m+M)$. Энергия второй шайбы $Mu^2/2=m^2Mv^2/[2(m+M)^2]$ (формула $3$); импульс первой $mu=m^2v/(m+M)$ (формула $1$). Ответ: $31$.
}

% Source F4CD91: https://ege.fipi.ru/bank/questions.php?search=1&proj=BA1F39653304A5B041B656915DC36B38&qid=F4CD91
\BankTask{F4CD91}{
Тело, брошенное с горизонтальной поверхности Земли со скоростью $v$ под углом $\alpha$ к горизонту, поднимается на максимальную высоту $h$ над земной поверхностью и затем падает на землю. Сопротивлением воздуха пренебречь.
\TaskMatchingColumns[.53\linewidth]{Физические величины}{
А) максимальная высота $h$ над горизонтом\par
Б) модуль вертикальной проекции скорости тела непосредственно перед падением на землю
}{Формулы}{
1) $\frac{v^2\sin^2\alpha}{2g}$\par
2) $v\sin\alpha$\par
3) $\frac{v^2\sin2\alpha}{2g}$\par
4) $v\cos\alpha$\par
}\begin{center}\begin{tabular}{|c|c|}\hline А & Б \\\hline\rule{0pt}{1.2em}\hspace{1cm} & \hspace{1cm} \\\hline\end{tabular}\end{center}
}{
\smallskip\noindent\textbf{Решение от GPT}\par
Начальная вертикальная скорость $v_{0y}=v\sin\alpha$. Из $v_y^2=v_{0y}^2-2gh$ при $v_y=0$ получаем $h=v^2\sin^2\alpha/(2g)$ (формула $1$). При возвращении на исходный уровень $v_y=-v\sin\alpha$, поэтому $|v_y|=v\sin\alpha$ (формула $2$). Ответ: $12$.
}

% Source B28492: https://ege.fipi.ru/bank/questions.php?search=1&proj=BA1F39653304A5B041B656915DC36B38&qid=B28492
\BankTask{B28492}{
Тело массой $200\,\text{г}$ движется вдоль оси $Ox$, при этом его координата изменяется во времени в соответствии с формулой $x(t)=10+5t-3t^2$ (все величины выражены в СИ).
\TaskMatchingColumns[.53\linewidth]{Физические величины}{
А) проекция $v_x(t)$ скорости тела\par
Б) проекция $F_x(t)$ равнодействующей сил, приложенных к телу
}{Формулы}{
1) $5-6t$\par
2) $-1{,}2$\par
3) $-3$\par
4) $10+5t$\par
}\begin{center}\begin{tabular}{|c|c|}\hline А & Б \\\hline\rule{0pt}{1.2em}\hspace{1cm} & \hspace{1cm} \\\hline\end{tabular}\end{center}
}{
\smallskip\noindent\textbf{Решение от GPT}\par
$v_x=dx/dt=5-6t$ (формула $1$), $a_x=-6\,\text{м}/\text{с}^2$. Тогда $F_x=ma_x=0{,}2\cdot(-6)=-1{,}2\,\text{Н}$ (формула $2$). Ответ: $12$.
}

% Source 69AD85: https://ege.fipi.ru/bank/questions.php?search=1&proj=BA1F39653304A5B041B656915DC36B38&qid=69AD85
\BankTask{69AD85}{
Выражения А и Б определяют зависимость координат двух тел от времени. Установите соответствие между зависимостью координаты тела от времени и зависимостью проекции скорости от времени для этого же тела. Все величины заданы в СИ.
\TaskMatchingColumns[.53\linewidth]{Координата}{
А) $x=10+5t+2t^2$\par
Б) $x=10+2t^2$
}{Скорость}{
1) $v_x=5+4t$\par
2) $v_x=5+2t$\par
3) $v_x=4t$\par
4) $v_x=2t$\par
}\begin{center}\begin{tabular}{|c|c|}\hline А & Б \\\hline\rule{0pt}{1.2em}\hspace{1cm} & \hspace{1cm} \\\hline\end{tabular}\end{center}
}{
\smallskip\noindent\textbf{Решение от GPT}\par
Скорость --- производная координаты: в случае А $v_x=5+4t$ (формула $1$), в случае Б $v_x=4t$ (формула $3$). Ответ: $13$.
}

% Source DB3508: https://ege.fipi.ru/bank/questions.php?search=1&proj=BA1F39653304A5B041B656915DC36B38&qid=DB3508
\BankTask{DB3508}{
Шайба массой $m$ съезжает без трения из состояния покоя с вершины горки. У подножия горки потенциальная энергия шайбы равна нулю, а модуль её импульса равен $p$. Каковы высота горки и потенциальная энергия шайбы на её вершине? Ускорение свободного падения равно $g$.
\TaskMatchingColumns[.53\linewidth]{Физические величины}{
А) высота горки\par
Б) потенциальная энергия шайбы на вершине горки
}{Формулы}{
1) $\frac{p^2}{2m^2g}$\par
2) $\frac{p^2}{2mg}$\par
3) $\frac{p^2}{2m}$\par
4) $\frac{mp^2}{2g}$\par
}\begin{center}\begin{tabular}{|c|c|}\hline А & Б \\\hline\rule{0pt}{1.2em}\hspace{1cm} & \hspace{1cm} \\\hline\end{tabular}\end{center}
}{
\smallskip\noindent\textbf{Решение от GPT}\par
$mgh=mv^2/2=p^2/(2m)$. Поэтому $h=p^2/(2m^2g)$ (формула $1$), $E_{\text{п}}=p^2/(2m)$ (формула $3$). Ответ: $13$.
}

% Source BE8E13: https://ege.fipi.ru/bank/questions.php?search=1&proj=BA1F39653304A5B041B656915DC36B38&qid=BE8E13
\BankTask{BE8E13}{
Один конец лёгкой пружины жёсткостью $k$ прикреплён к бруску, а другой закреплён неподвижно. Брусок скользит по горизонтальной направляющей так, что его координата изменяется со временем по закону $x(t)=A\sin\omega t$.
Установите соответствие между физическими величинами и формулами, выражающими их изменения во времени.\TaskMatchingColumns[.53\linewidth]{Физические величины}{
А) проекция $F_x(t)$ равнодействующей силы на ось $x$\par
Б) потенциальная энергия пружины $E_{\text{П}}(t)$
}{Формулы}{
1) $kA^2\sin\omega t$\par
2) $-kA\sin\omega t$\par
3) $-kA^2\sin^2\omega t$\par
4) $\frac{kA^2}{2}\sin^2\omega t$\par
}\begin{center}\begin{tabular}{|c|c|}\hline А & Б \\\hline\rule{0pt}{1.2em}\hspace{1cm} & \hspace{1cm} \\\hline\end{tabular}\end{center}
}{
\smallskip\noindent\textbf{Решение от GPT}\par
По закону Гука $F_x=-kx=-kA\sin\omega t$ (формула $2$). Энергия $E_{\text{П}}=kx^2/2=(kA^2/2)\sin^2\omega t$ (формула $4$). Ответ: $24$.
}

% Source DF724F: https://ege.fipi.ru/bank/questions.php?search=1&proj=BA1F39653304A5B041B656915DC36B38&qid=DF724F
\BankTask{DF724F}{
Подвешенный на пружине груз совершает свободные вертикальные гармонические колебания. Пружину заменили на другую, жёсткость которой больше, оставив массу груза и амплитуду колебаний неизменными. Как при этом изменятся частота свободных колебаний груза и его максимальная скорость?
\begin{enumerate}
\item увеличится
\item уменьшится
\item не изменится
\end{enumerate}
\TaskChangeTable{Частота свободных колебаний груза}{Максимальная скорость груза}
}{
\smallskip\noindent\textbf{Решение от GPT}\par
$\nu=(2\pi)^{-1}\sqrt{k/m}$, $v_{\max}=2\pi\nu A$. При увеличении $k$ и неизменных $m,A$ обе величины увеличиваются. Ответ: $11$.
}

% Source F4D3FF: https://ege.fipi.ru/bank/questions.php?search=1&proj=BA1F39653304A5B041B656915DC36B38&qid=F4D3FF
\BankTask{F4D3FF}{
Космический аппарат, обращающийся вокруг Луны по круговой орбите, перешёл на другую круговую орбиту бóльшего радиуса. Как изменились в результате этого перехода центростремительное ускорение, с которым аппарат движется по орбите, и его период обращения вокруг Луны?
\begin{enumerate}
\item увеличилась
\item уменьшилась
\item не изменилась
\end{enumerate}
\TaskChangeTable{Центростремительное ускорение}{Период обращения аппарата вокруг Луны}
}{
\smallskip\noindent\textbf{Решение от GPT}\par
$a=GM/r^2$ уменьшается с ростом радиуса; $T=2\pi\sqrt{r^3/(GM)}$ увеличивается. Ответ: $21$.
}

% Source FcA2F2: https://ege.fipi.ru/bank/questions.php?search=1&proj=BA1F39653304A5B041B656915DC36B38&qid=FcA2F2
\BankTask{FcA2F2}{
В школьном опыте брусок, помещённый на горизонтальный диск, вращался вместе с ним с некоторой угловой скоростью. В ходе опыта угловую скорость диска увеличили. При этом положение бруска на диске осталось прежним. Как изменились при этом центростремительное ускорение бруска и сила нормального давления бруска на опору?
\begin{enumerate}
\item увеличилась
\item уменьшилась
\item не изменилась
\end{enumerate}
\TaskChangeTable{Центростремительное ускорение бруска}{Сила нормального давления бруска на опору}
}{
\smallskip\noindent\textbf{Решение от GPT}\par
$a=\omega^2r$ увеличивается. По вертикали ускорения нет, поэтому сила нормального давления $N=mg$ не меняется. Ответ: $13$.
}

% Source 589e07: https://ege.fipi.ru/bank/questions.php?search=1&proj=BA1F39653304A5B041B656915DC36B38&qid=589e07
\BankTask{589e07}{
Мяч подбросили вверх с крыши сарая. Как изменяются по мере подъёма модуль импульса мяча и его потенциальная энергия относительно поверхности Земли? Сопротивление воздуха не учитывать.
\begin{enumerate}
\item увеличивается
\item уменьшается
\item не изменяется
\end{enumerate}
\TaskChangeTable{Модуль импульса мяча}{Потенциальная энергия мяча}
}{
\smallskip\noindent\textbf{Решение от GPT}\par
При подъёме скорость уменьшается, поэтому $p=mv$ уменьшается. Высота растёт, значит $E_{\text{п}}=mgh$ увеличивается. Ответ: $21$.
}

% Source CCDA00: https://ege.fipi.ru/bank/questions.php?search=1&proj=BA1F39653304A5B041B656915DC36B38&qid=CCDA00
\BankTask{CCDA00}{
На поверхности воды плавает прямоугольный брусок из древесины плотностью $400\,\text{кг}/\text{м}^{3}$. Брусок заменили на другой брусок той же массы и с той же площадью основания, но из древесины плотностью $600\,\text{кг}/\text{м}^{3}$. Как при этом изменились глубина погружения бруска и действующая на него сила Архимеда?
\begin{enumerate}
\item увеличилась
\item уменьшилась
\item не изменилась
\end{enumerate}
\TaskChangeTable{Глубина погружения бруска}{Сила Архимеда}
}{
\smallskip\noindent\textbf{Решение от GPT}\par
При плавании $F_A=mg$, а $\rho_{\text{воды}}gSh=mg$. При одинаковых $m,S$ обе величины неизменны: $h=m/(\rho_{\text{воды}}S)$. Ответ: $33$.
}

% Source AEDF70: https://ege.fipi.ru/bank/questions.php?search=1&proj=BA1F39653304A5B041B656915DC36B38&qid=AEDF70
\BankTask{AEDF70}{
Железный сплошной шарик совершает малые свободные колебания на лёгкой нерастяжимой нити. Затем этот шарик заменили на сплошной алюминиевый шарик такого же диаметра. Амплитуда колебаний в обоих случаях одинакова.
Как при этом изменятся частота колебаний и максимальная кинетическая энергия шарика?
\begin{enumerate}
\item увеличится
\item уменьшится
\item не изменится
\end{enumerate}
\TaskChangeTable{Частота колебаний шарика}{Максимальная кинетическая энергия шарика}
}{
\smallskip\noindent\textbf{Решение от GPT}\par
Длина маятника одинакова, поэтому $\nu=(2\pi)^{-1}\sqrt{g/l}$ не меняется. При одинаковой амплитуде $v_{\max}$ одинакова, а масса алюминиевого шарика меньше, поэтому $E_{\text{к max}}=mv_{\max}^2/2$ уменьшается. Ответ: $32$.
}

% Source 010677: https://ege.fipi.ru/bank/questions.php?search=1&proj=BA1F39653304A5B041B656915DC36B38&qid=010677
\BankTask{010677}{
Камень подбросили вверх. Как меняются по мере подъёма ускорение камня
и его кинетическая энергия? Сопротивление воздуха не учитывать.
\begin{enumerate}
\item увеличивается
\item уменьшается
\item не меняется
\end{enumerate}
\TaskChangeTable{Ускорение камня}{Кинетическая энергия камня}
}{
\smallskip\noindent\textbf{Решение от GPT}\par
Ускорение свободного падения $\vec a=\vec g$ постоянно. При подъёме скорость уменьшается, поэтому $E_{\text{к}}=mv^2/2$ уменьшается. Ответ: $32$.
}

% Source E2791B: https://ege.fipi.ru/bank/questions.php?search=1&proj=BA1F39653304A5B041B656915DC36B38&qid=E2791B
\BankTask{E2791B}{
На поверхности подсолнечного масла плавает деревянный шарик, частично погружённый в жидкость. Как изменятся сила Архимеда, действующая
на шарик, и глубина погружения шарика, если он будет плавать в воде?
\begin{enumerate}
\item увеличится
\item уменьшится
\item не изменится
\end{enumerate}
\TaskChangeTable{Сила Архимеда}{Глубина погружения шарика в жидкость}
}{
\smallskip\noindent\textbf{Решение от GPT}\par
При плавании $F_A=mg$ не меняется. Плотность воды больше плотности масла, поэтому $V_{\text{погр}}=m/\rho$ и глубина погружения уменьшаются. Ответ: $32$.
}

% Source 32F41D: https://ege.fipi.ru/bank/questions.php?search=1&proj=BA1F39653304A5B041B656915DC36B38&qid=32F41D
\BankTask{32F41D}{
Шарик, брошенный с поверхности Земли под углом $20\,{}^\circ$ к горизонту
с начальной скоростью ${\vec{\upsilon}}_0$, поднялся на максимальную высоту $H$ и пролетел
в горизонтальном направлении расстояние $L$. Что произойдёт с дальностью полёта шарика и временем его полёта, если шарик бросить с той же
по величине начальной скоростью под углом $30\,{}^\circ$ к горизонту?
Сопротивлением воздуха пренебречь.
\begin{enumerate}
\item увеличится
\item уменьшится
\item не изменится
\end{enumerate}
\TaskChangeTable{Дальность полёта шарика}{Время полёта шарика}
}{
\smallskip\noindent\textbf{Решение от GPT}\par
$L=v_0^2\sin2\alpha/g$, $t=2v_0\sin\alpha/g$. При росте угла от $20^\circ$ до $30^\circ$ возрастают и $\sin2\alpha$, и $\sin\alpha$, следовательно, обе величины увеличиваются. Ответ: $11$.
}

% Source FB0F2e: https://ege.fipi.ru/bank/questions.php?search=1&proj=BA1F39653304A5B041B656915DC36B38&qid=FB0F2e
\BankTask{FB0F2e}{
В первой серии опытов брусок перемещают при помощи нити равномерно вверх по наклонной плоскости. Во второй серии опытов на брусок ставят груз и продолжают перемещать брусок при помощи нити равномерно вверх по той же наклонной плоскости. В обоих случаях нить параллельна наклонной плоскости.
Как изменяются при переходе от первой серии опытов ко второй модуль работы силы трения при перемещении бруска на одинаковые расстояния и коэффициент трения скольжения между бруском и плоскостью?
\begin{enumerate}
\item увеличивается
\item уменьшается
\item не изменяется
\end{enumerate}
\TaskChangeTable{Модуль работы силы трения}{Коэффициент трения скольжения}
}{
\smallskip\noindent\textbf{Решение от GPT}\par
При той же плоскости $\mu$ неизменен. После добавления груза $N=(m+M)g\cos\alpha$ больше, поэтому $|A_{\text{тр}}|=\mu Ns$ на том же пути увеличивается. Ответ: $13$.
}

% Source 10BE25: https://ege.fipi.ru/bank/questions.php?search=1&proj=BA1F39653304A5B041B656915DC36B38&qid=10BE25
\BankTask{10BE25}{
На поверхности пресной воды плотностью $\rho_{1}=1000\,\text{кг}/\text{м}^{3}$ плавает деревянный брусок. Как изменятся масса вытесненной бруском жидкости и действующая на него сила Архимеда, если этот брусок будет плавать на поверхности керосина плотностью $\rho_{2}=800\,\text{кг}/\text{м}^{3}$?
\begin{enumerate}
\item увеличится
\item уменьшится
\item не изменится
\end{enumerate}
\TaskChangeTable{Масса вытесненной бруском жидкости}{Сила Архимеда}
}{
\smallskip\noindent\textbf{Решение от GPT}\par
При плавании $F_A=mg$, а масса вытесненной жидкости $m_{\text{ж}}=F_A/g=m$. Обе величины не зависят от плотности жидкости и не изменяются. Ответ: $33$.
}

% Source A5cD6A: https://ege.fipi.ru/bank/questions.php?search=1&proj=BA1F39653304A5B041B656915DC36B38&qid=A5cD6A
\BankTask{A5cD6A}{
Шарик, брошенный с поверхности Земли под углом $\alpha$ к горизонту с начальной скоростью ${\vec{\upsilon}}_0$, поднялся на максимальную высоту $H$ и пролетел в горизонтальном направлении расстояние $L$. Что произойдёт с максимальной высотой, на которую поднимется шарик, и его ускорением, если уменьшить начальную скорость шарика, не меняя $\alpha$? Сопротивлением воздуха пренебречь.
\begin{enumerate}
\item увеличится
\item уменьшится
\item не изменится
\end{enumerate}
\TaskChangeTable{Максимальная высота подъёма}{Ускорение шарика}
}{
\smallskip\noindent\textbf{Решение от GPT}\par
$H=v_0^2\sin^2\alpha/(2g)$ уменьшается при уменьшении $v_0$ с неизменным углом. Ускорение $\vec a=\vec g$ не меняется. Ответ: $23$.
}

% Source 9A9FD1: https://ege.fipi.ru/bank/questions.php?search=1&proj=BA1F39653304A5B041B656915DC36B38&qid=9A9FD1
\BankTask{9A9FD1}{
Спортсмен спускается на парашюте с постоянной скоростью. Как изменяются с течением времени в процессе спуска импульс спортсмена и его потенциальная энергия?
\begin{enumerate}
\item увеличивается
\item уменьшается
\item не изменяется
\end{enumerate}
\TaskChangeTable{Импульс}{Потенциальная энергия}
}{
\smallskip\noindent\textbf{Решение от GPT}\par
При постоянной скорости $\vec p=m\vec v$ не меняется. Высота уменьшается, поэтому $E_{\text{п}}=mgh$ уменьшается. Ответ: $32$.
}

\end{document}
